\section{classification}

Differential equations are a class of \emph{functional equations}
\mymargin{functional equations}
\index{equation!functional}
that is, equations in which the unknown is not a number (as in algebraic
equations) but a function. The unknown function $u$ will therefore be a function
of an independent variable $u=u(x)$.
For example, $u$ could be the trajectory of a projectile that describes the
position in space as a function of time $x$. While in algebraic equations the
most commonly used name for the unknown is $x$, in functional equations,
depending on the context, conventions can change drastically. One must use a
name for the function and a name for its independent variable: $x=x(t)$,
$y=y(x)$, $y=y(t)$, $u=u(x)$ are some of the most commonly used choices.
If the unknown is a function $u=u(x)$, the operations that can appear in the
equations are, in addition to the usual algebraic operations that act on the
individual values $u(x)$ of the function, also operators that act on the
function $u$ itself.
If the functional equation, in addition to algebraic operations, also includes
the derivative operator, it is called a \mymargin{differential equation}%
\index{equation!differential}%
\emph{differential equation}.
Conversely, there may be equations that involve the integral operator, which are
called \emph{integral equations}.

There are several concepts that can help classify differential equations.
First of all, we will always assume that our equations are \emph{without delay}%
\mymargin{without delay}\index{without delay}
(or \emph{without memory}),
meaning that the function $u$ and its derivatives $u'$, $u''$\dots,
are all calculated at the same point $x$.
Otherwise, we speak of
\emph{functional differential equations} (in particular \emph{equations with delay}
because the derivatives typically depend on values assumed in the past),
a topic we will not address.

If the unknown function is a function of a single variable, the equation is
called an
\emph{ordinary differential equation}
\mymargin{ODE}%
\index{equation!differential!ordinary}%
(abbreviated
\emph{ODE} \index{ODE}%
in English).
Conversely, if the unknown function is a function of multiple variables, the
equation is called a
\emph{partial differential equation}
\mymargin{PDE}%
\index{equation!differential!partial}%
(abbreviated \emph{PDE}
\index{PDE} in English).
In this course, focused on functions of a single variable, we will therefore
only deal with ordinary differential equations.
The \emph{order}
\mymargin{order}%
\index{order!differential equation}%
\index{equation!differential!order}%
of the differential equation is the maximum number of successive derivatives
applied to the unknown function.
The function $u$ will be defined on an interval $I$ of the real line: $u\colon I
\to \RR$
and the function and its derivatives are all evaluated at the same point. In
this case,
the most general form of an ordinary differential equation of order $n$ can be
written as:
\begin{equation}\label{eq:3784643}
  F(x,u(x),u'(x),u''(x), \dots, u^{(n)}(x)) = 0
\end{equation}
with $F\colon \Omega \to \RR$ a given function, defined on a set $\Omega \subset
I\times \RR^{n+1}$.
A function $u\colon I \to \RR$ is said to be a \emph{solution
of the differential equation} \eqref{eq:3784643} if $u$ is differentiable at
least
$n$ times at every point $x\in I$ and if \eqref{eq:3784643} is satisfied
for every $x\in I$.

Usually, the notation is simplified by avoiding writing explicitly
the point $x$ at which the function is calculated.
It will therefore be usual
to write the equation \eqref{eq:3784643}
in the abbreviated form:
\[
  F(x,u,u',u'', \dots, u^{(n)}) = 0
\]
making it more evident that the unknown is $u$,
the entire function, and not a single value $u(x)$.

For example, if we choose $n=2$ and $F(x,u,v,z)= z+\sin u + v$, we obtain the
differential equation:
\begin{equation}\label{eq:39872}
 u''(x) + \sin u(x) + u'(x) = 0
\end{equation}
which, in appropriate units of measurement, is the equation of motion of a
damped pendulum,
where $x$ represents time and $u$ the measure of the angle of inclination of the
pendulum with respect to the vertical.
We observe that in the previous example the function $F$ does not directly
depend
on the variable $x$. Equations with this property are called
\emph{autonomous equations}
\mymargin{autonomous equations}
\index{equation!differential!autonomous}
and it is immediate to observe that if $u(x)$ is a solution, then a temporal
translation
$v(x) = u(x-x_0)$ is also a solution of the equation
(the motion of the pendulum does not depend on the time at which it occurs).

An equation written in the form \eqref{eq:3784643} is called an \emph{implicit
form equation}
\mymargin{implicit form equation}%
\index{equation!differential!implicit form}%
and by analogy with algebraic equations (consider the equation
$u^2(x) + x^2 = 1$) it is expected that the solutions of such equations are
better represented by curves rather than function graphs.
It turns out that the theory of differential equations applies with
much greater effectiveness to \emph{normal form equations}
\mymargin{normal form equations}
\index{equation!differential!normal form}
which are differential equations of order $n$ that can be written
by making explicit the dependence on the highest order derivative:
\begin{equation}\label{eq:366793}
 u^{(n)}(x) = f(x, u(x), u'(x), \dots, u^{n-1}(x))
\end{equation}
where $f\colon \Omega \to \RR$ is a function defined on $\Omega\subset I\times
\RR^n$.
The pendulum equation can be written in this form,
choosing $f(x,u,v) = -\sin u - v$.

More generally, we could consider
\emph{systems of differential equations}
\mymargin{systems}%
\index{systems of differential equations}%
\index{equation!differential!system}%
in multiple unknowns.
We can represent a system of $k$ ordinary equations in $m$ unknowns in the form:
\begin{equation}\label{eq:375456}
  \vec F(x, \vec u(x), \vec u'(x), \dots, \vec u^{n}(x)) = \vec 0
\end{equation}
where $\vec u$ is a vector function $\vec u \colon I \to \RR^m$ 
whose components are the $m$ unknown functions:
\[
  \vec u(x) = (u_1(x), \dots, u_m(x))
\]
while the function $\vec F \colon I \times (\RR^m)^{n+1} \to \RR^k$ is now
a vector-valued function $\vec F = (F_1, \dots, F_k)$ so that the vector
equation
\eqref{eq:375456} is effectively equivalent to a system of $k$ equations:
\[
\begin{cases}
F_1(x,\vec u(x), \vec u'(x)\dots, \vec u^{(n)}(x)) = 0\\
F_2(x,\vec u(x), \vec u'(x)\dots, \vec u^{(n)}(x)) = 0\\
\quad\vdots \\
F_k(x,\vec u(x), \vec u'(x)\dots, \vec u^{(n)}(x)) = 0\\
\end{cases}
\]
We will see that for first-order ordinary equations, it is natural,
as happens for algebraic equations,
to have the same number of equations and unknowns, thus it is typical to have
single
scalar first-order equations
(that is, where the unknown is a function valued in the scalar field
$\RR$)
or systems of $k=m$ first-order equations with the unknown being a vector
function
$\vec u$ (that is, a function valued in the vector space $\RR^m$)
or with $m$ scalar unknowns $u_1,\dots, u_m$.

An important observation is the general fact that a differential equation
of order $n$ can be reduced to a system of $n$ first-order differential
equations.
Indeed, it suffices to consider as the unknown the vector (called \emph{jet})
\[
  \vec u = (u, u', u'', \dots, u^{(n-1)})
\]
comprising all derivatives of the scalar function $u$ up to order $n-1$.
The equation \eqref{eq:3784643}, of order $n$, is indeed equivalent
to the system of $n$ first-order equations in the variable
$\vec u = (u_1, \dots, u_n)$
\[
  \begin{cases}
    u_2(x) = u_1'(x)\\
    u_3(x) = u_2'(x)\\
    \quad \vdots \\
    u_{n}(x) = u_{n-1}'(x)\\
    F(x, u_1(x), u_2(x), \dots, u_n(x), u_n'(x)) = 0
  \end{cases}
\]
In the more interesting case,
of normal form equations,
the equation~\eqref{eq:366793} of order $n$
becomes a system of $n$ normal first-order equations
in $n$ unknowns
\[
  \begin{cases}
  u_1'(x) = u_2(x)\\
  u_2'(x) = u_3(x)\\
  \quad \vdots\\
  u_{n-1}'(x) = u_n(x)\\
  u_n'(x) = f(x,u_1(x), u_2(x), \dots, u_n(x))
  \end{cases}
\]
that is, a first-order vector differential equation
\[
  \vec u'(x) = \vec f(x, \vec u(x)).
\]
\begin{comment}
having defined $\vec f = (f_1, \dots, f_n)$
as
\begin{gather*}
  f_1(x, y_1,\dots, y_n) = y_2 \\
  f_1(x, y_1,\dots, y_n) = y_3 \\
  \quad \vdots \\
  f_{n-1}(x, y_1,\dots, y_n) = y_n \\
  f_n(x, y_1,\dots, y_n) = f(x,y_1, \dots, y_n)
\end{gather*}
\end{comment}

In the example of the pendulum \eqref{eq:39872}, the unknown will be a vector
function
$\vec u(x) = (u(x),u'(x))$ whose components are position and angular velocity.
The codomain of this function is called the \emph{phase space}.
The equation (being autonomous, we omit the dependence on $x$)
will be written in the form $\vec u' = \vec f(\vec u)$ with
$\vec f(y_1,y_2) = (y_2, -\sin(y_1) - y_2)$.

\subsection{linear equations}

A very particular but decidedly important case is when the function
$F$ (for implicit form equations) or the function $f$
(for normal form equations) are linear functions
for each $t$
with respect to the variable $u$.
In this case, we will say that the equation is
\mymargin{homogeneous linear equations}%
\index{equation!differential!homogeneous linear}%
\emph{homogeneous linear}.
More precisely, we will have\mynote{%
To highlight the fact that a function $L$ is linear, we will tend to use 
square brackets (instead of parentheses) to enclose the argument of the 
function: $L[v] = L(v)$.
}
\[
F(x,u(x),u'(x), \dots, u^{(n)}(x)) = A_x[u(x), u'(x), \dots, u^{(n)}(x)]
\]
with $A_x\colon \RR^{n+1}\to \RR$
a linear operator for each $x$, that is, $A_x$ is represented by a vector
whose coefficients are functions of the variable $x$:
\[
  A_x[\vec y] = \sum_{k=0}^n a_k(x) y_k
\]
and the differential equation is written in the form
\[
  a_0(x) u(x) + a_1(x) u'(x) + \dots + a_n(x) u^{(n)}(x) = 0.
\]
In the case where the coefficients $a_k(t)$ do not depend on $x$
(that is, they are constant functions), we will say that the equation is
linear
\mymargin{constant coefficients}%
\index{equation!differential!linear with constant coefficients}%
\emph{with constant coefficients}.

It is easy to observe that the set of solutions of a homogeneous linear equation
is a vector space: $u=0$ is always a solution,
if $u$ is a solution and $\lambda \in \RR$, then $\lambda u$ is also a solution,
and if $u$ and $v$ are two solutions, then $u+v$ is a solution
\mymargin{superposition principle}%
\index{principle!of superposition}%
\index{principle!of superposition}%
(superposition principle).

Indeed,
if the functions $u$ are defined on an interval $I$, we can
identify $F$ with a functional $L\colon \RR^I \to \RR^I$ defined by
\[
  L[u](x) = F(x,u(x), u'(x), \dots, u^{(n)}(x))
\]
If the differential equation is linear, then $L$ is a linear operator
on the vector space $\RR^I$, and the space of solutions of the
differential equation is nothing but $\ker L$, the kernel of the operator,
and it is known that $\ker L$ is a vector subspace.
Although $\RR^I$ is a vector space of infinite dimension, we will discover
that (under appropriate assumptions) the vector space of solutions
has finite dimension $n$, equal to the degree of the equation.

In the case where the function $F$ (or the corresponding $f$)
is affine\mynote{%
A function $f\colon V\to W$ defined between two vector spaces is said 
to be \emph{affine} if it can be written in the form $f(v)=L[v] + q$
with $L\colon V\to W$ a linear operator and $q\in W$ fixed.
A linear equation $L[v]=w$ can therefore always be written in the form 
$f(v)=0$ with $f$ affine.
} 
the differential equation
is called a
\mymargin{linear equation}%
\index{equation!differential!linear}%
\emph{linear (non-homogeneous) equation}.
The non-homogeneous equation will have the form:
\[
  a_0(x) u(x) + a_1(x) u'(x) + \dots + a_n(x) u^{(n)}(x) = g(x).
\]
If $v_0$ and $v_1$ are two solutions of this equation, it is clear that the
difference $u=v_1 - v_0$ is a solution of the homogeneous equation
\[
  a_0(x) u(x) + a_1(x) u'(x) + \dots + a_n(x) u^{(n)}(x) = 0.
\]
Thus, the latter is called the homogeneous equation associated with the
non-homogeneous one,
and if $v_0$ is a particular (any) solution of the non-homogeneous equation,
every solution $v$ of the non-homogeneous equation can be written in the form
\[
  v = v_0 + u
\]
with $u$ a solution of the associated homogeneous equation.
To find all solutions of a non-homogeneous equation, it is therefore
sufficient to find a particular solution of the non-homogeneous equation 
and all solutions of the associated homogeneous equation.

The pendulum equation \eqref{eq:39872} is not linear,
but when the angle $u$ is small (i.e., for \emph{small oscillations}), we have
$\sin u \sim u$.
By making this \emph{linearization}, we obtain the equation
\[
  u''(x)  + u(x) + u'(x) = 0.
\]
This is a homogeneous linear equation.
If an external force acts on the pendulum (forced pendulum), the equation
becomes
\[
  u''(x) + u(x) + u'(x) = g(x)
\]
where $g(x)$ represents the magnitude of an external force varying over time.
This equation is a non-homogeneous linear equation.

As in the general case, linear equations of order $n$ can be reduced
to linear systems of $n$ first-order equations.

\subsection{Cauchy problem}

We will demonstrate not only that linear equations of order $n$ (homogeneous or
non-homogeneous)
have as their set 
of solutions a linear (vector or affine) space of dimension $n$, 
but, more generally, that the set of solutions of a differential equation, 
even if not linear, can be described by $n$ parameters
(therefore, this set of solutions can be thought of as a ``surface'' 
of dimension $n$ in the infinite-dimensional vector space 
of all differentiable functions).

The idea is that ordinary equations of order $n$ 
have a form of \emph{determinism}:
if I know the solution $u(x_0)$ at a moment $x_0$ 
and I also know all its derivatives $u'(x_0), u''(x_0), \dots, u^{(n-1)}(x_0)$ 
at that moment, then the solution exists and is uniquely determined
(Theorem~\ref{th:cauchy_lipschitz_ordine_n}).

The values $u(x_0), u'(x_0), \dots, u^{(n-1)}(x_0)$ are called \emph{initial
conditions}.
and the problem of determining a solution of a differential equation
\mymargin{initial condition}%
\index{initial condition}%
with fixed initial conditions is called the \emph{Cauchy problem}:
\mymargin{Cauchy problem}%
\index{problem!Cauchy}%
\[
  \begin{cases}
  u^{(n)}(x) = f(x, u(x), \dots, u^{n-1}(x))\\
  u(x_0) = y_0, \\
  u'(x_0) = y_1, \\
  \vdots\\
  u^{(n-1)}(x_0) = y_{n-1}.
  \end{cases}
\]
We will see that, under appropriate assumptions, the Cauchy problem
admits a unique solution, and this means, essentially,
that the set of all solutions of the differential equation
can be described by varying
the $n$ parameters $y_0, y_1, \dots, y_{n-1}$.

One of the fundamental assumptions to guarantee the uniqueness of
solutions to a Cauchy problem is that the solutions are
defined on an interval. It is clear that if I have two
solutions defined on two separate intervals, I could join
the two solutions and obtain a solution defined on the union of
the two intervals while being able to choose independent initial conditions 
in each of the two intervals. 
Conversely, a solution defined
on an interval $I$ is also a solution when restricted to any subset
of $I$ (whether or not it is an interval).
If I restrict a solution or join two solutions, I formally obtain
different solutions since the domains of definition are different, but
these solutions assume the same values where they are defined.
To avoid this ambiguity, we will always consider only
solutions defined on a \emph{maximal interval}%
\mymargin{maximal interval}\index{maximal interval} that is,
solutions defined on an interval that cannot be
extended to larger intervals.